\section{A single standard layer: exact descent and the extra summand}
\label{sec:standard-single-factor}

\paragraph{Status and scope.}
HD-SS01--03, reserved question Q2ea, is a separate derivation by
\texttt{/root/imprimitive}.  The parent proposed retaining the exact
induction identity and testing whether its extra cohomology vanishes.
The restriction and character-field argument below is new in this
derivation.  The counterexample concerns an Euler-only inference;
it has rational coefficient field and period rank at most six.
It is not a counterexample for an exact imaginary-quadratic, rank-nine
sixfold supplier.  The latter remaining question is stated explicitly.
The complete author proof, source-interface and certificate readings
pass. Complete root and independent whole-proof, source and code
readings also pass without mathematical repair. No main-target exclusion
is asserted by the counterexample.

Let $J$ be an actual finite coefficient image, and suppose its faithful
absolutely irreducible complex representation has a factorization
\[
 V=W\otimes U,\qquad W=\operatorname{Std}_n\circ x,
 \qquad x:J\longrightarrow S_n,
 \qquad A_n\subset x(J),\quad n\ge4.
 \tag{SS.1}
\]
In the stable-layer applications here $n\ge15$, and its compatible
projective normalizer has the unique rational lift to $S_n$ used
in HD-AST13--15.  Division by that lift gives an actual homomorphism
$U:J\to\operatorname{GL}_b(\mathbf C)$, not just a projective
representation; it has finite image because $J$ is finite.
For the smaller $n$ allowed in (SS.1), the displayed actual
factorization is the hypothesis, and no additional normalizer
classification is being asserted.  Thus (SS.1) also applies to
a deeper layer of the recursive
decomposition D-TD11--13.  It does not assume that the projected
$A_n$ lifts to a subgroup acting trivially on $U$.
Put $a=n-1$, so $\dim V=ab$, and let
$H=x^{-1}((S_n)_1)$ be the stabilizer of one letter.
We write $F_T=\mathbf Q(\chi_T)$ for an exact character field.

\subsection*{HD-SS01. Irreducible restriction and exact fields}

\begin{proposition}[HD-SS01]
The representation $U$ and its restriction $\tau=\operatorname{Res}_H^J U$
are irreducible.  There is an actual representation identity
\[
 \operatorname{Ind}_H^J\tau=U\oplus V,
 \qquad F_\tau=F_U\supset F_V.
 \tag{SS.2}
\]
If $V$ has a model over its exact imaginary quadratic field $K=F_V$,
put $L=F_U=F_\tau$.  Both $U$ and $\tau$ have models over $L$, their
Schur indices over $L$ are one, and
\[
                         [L:K]\le n-1.
 \tag{SS.3}
\]
In particular $U$ is nonselfdual in this exact-$K$ situation.
No equality $U(H)=U(J)$ is asserted.
\end{proposition}

\begin{proof}
A proper $J$-invariant subspace of $U$ would give a proper invariant
subspace after tensoring with $W$, so $U$ is irreducible.
The actual transitive letter permutation module is
$\operatorname{Ind}_H^J1=1\oplus W$.
The projection formula, in its explicit tensor-induction form, gives
the first identity of (SS.2).  Frobenius reciprocity then yields
\[
 \dim\operatorname{End}_H(\tau)
   =\dim\operatorname{Hom}_J(\operatorname{Ind}_H^J\tau,U)
   =\dim\operatorname{Hom}_J(U\oplus V,U)=1.
 \tag{SS.4}
\]
Here $\dim V=(n-1)b>b$, so the two irreducible constituents are distinct.
Semisimplicity in characteristic zero proves that $\tau$ is
irreducible.  The exact adjunction and projection formula are
Serre's Section~7.1, Remark~2, and Section~7.2, Theorem~13 and
Remark~3, including their proofs, pp.~55--57
\cite{OG:Serre1977}.

Restriction gives $F_\tau\subset F_U$.  Take a cyclotomic splitting
field containing every character value.  An automorphism fixing
$F_\tau$ fixes the induced character $\chi_U+\chi_V$.
Galois automorphisms permute irreducible characters and preserve
their degrees.  The two different degrees therefore force that
automorphism to fix each character separately.  Consequently
$F_U,F_V\subset F_\tau$, proving the field assertion.  Also,
because $W$ is rational, $\chi_V=\chi_W\chi_U$ directly gives
$F_V\subset F_U$.  A selfdual finite-group character is real;
it could not contain the imaginary field $K$ in its character field.

For the model assertion, Frobenius reciprocity applied to (SS.2)
shows that $\tau$ occurs exactly once in $\operatorname{Res}_H^J V$.
Its absolute central idempotent
\[
 e_\tau=\frac{b}{|H|}\sum_{h\in H}\chi_\tau(h^{-1})h
                 \ \in L[H]
 \tag{SS.5}
\]
applied to the $L$-module $V_K\otimes_KL$ has complexification
exactly $\tau$.  This constructs an $L$-model rather than assuming
that character-field equality implies realizability.
Induce this model to $J$ and apply the central idempotent of $V$,
which is already defined over $K$.  Its complement gives an
$L$-model of $U$.

Finally all $\ell=[L:K]$ distinct conjugates $\tau^\gamma$, where
$\gamma$ fixes $K$, occur once in $\operatorname{Res}_H V$.
Indeed conjugate the preceding multiplicity-one statement and use
the $K$-model of $V$.  Their total dimension $\ell b$ is at most
$(n-1)b$, proving (SS.3).
\end{proof}

\paragraph{The stronger pure-core interface has an extra hypothesis.}
If an actual normal subgroup mapping onto $A_n$ acts trivially on
$U$, then $J=NH$, so $U(H)=U(J)$.  In this case $F_U=F_V$:
for each $g\in J$, multiply by an element of $N$ so that its
permutation becomes the identity or a transposition according
to its parity.  Its $U$-matrix is unchanged, whereas the standard
trace becomes $n-1$ or $n-3$, both nonzero.  Thus
$\chi_U(g)\in F_V$.  None of these same-image or field-equality
conclusions is used without the indicated actual subgroup.

\subsection*{HD-SS02. Hodge and local Euler identities}

Let $C\to\mathbf P^1$ be a connected $J$-Galois cover with an
ordered generating product-one tuple $g_1,\ldots,g_r$.
Let $X=C/H$, $h=g(X)$, and put $N=\ker\tau\lhd H$ and $Y=C/N$.
The cover $Y\to X$ is connected, has faithful group $H/N$, and
has the irreducible coefficient representation $\tau$ of degree $b$.
In the statements concerning a sixfold, assume
\[
 \dim\operatorname{Hom}_J(V,H^1(C,\mathbf C))=6,
 \qquad (h^{1,0}_V,h^{0,1}_V)=(3,3),
 \qquad \mu_U=\dim\operatorname{Hom}_J(U,H^1(C,\mathbf C)).
 \tag{SS.6}
\]

\begin{proposition}[HD-SS02]
There is an isomorphism of complex weight-one multiplicity
variations, compatible with all coefficient embeddings,
\[
 \operatorname{Hom}_{H/N}(\tau,H^1(Y,\mathbf C))
 \simeq
 \operatorname{Hom}_J(U,H^1(C,\mathbf C))
 \oplus\operatorname{Hom}_J(V,H^1(C,\mathbf C)).
 \tag{SS.7}
\]
It is supplied by fixed group-algebra maps and quotient pullback.
If $x_i=x(g_i)$ and $S_i=U(g_i)$, then, with $c$ running over
the cycles of $x_i$,
\[
 \operatorname{rank}(1-V(g_i))
  =b\,\operatorname{ind}(x_i)
     +\sum_c\operatorname{rank}(1-S_i^{|c|})
     -\operatorname{rank}(1-S_i).
 \tag{SS.8}
\]
For nontrivial $U$, the total moving rank $D'$ of $\tau$ on
$Y\to X$, counting only genuine inertia, satisfies
\[
               6+\mu_U=2b(h-1)+D',
 \qquad D'=\sum_i\sum_c\operatorname{rank}(1-S_i^{|c|}).
 \tag{SS.9}
\]
For trivial $U$ the right side in the first equality has the
additional term $2$; the nontrivial exact-$K$ case always uses
(SS.9) as displayed.
\end{proposition}

\begin{proof}
Apply the adjunction underlying (SS.2) to $H^1(C)$.
The finite quotient pullback identifies $H^1(Y)$ with
$H^1(C)^N$.  Every map is a fixed linear combination of deck
correspondences or a finite quotient pullback, so it preserves
the Hodge filtration in families.  This proves (SS.7) over a
common splitting field and after every embedding.  The assertion
does not replace the sum by its second summand.

On a cycle of length $l$, the fixed vectors of
$\mathbf C^n(x_i)\otimes S_i$ are obtained by choosing a vector
fixed by $S_i^l$ and transporting it around that cycle.
Since $\mathbf C^n=1\oplus W$, subtracting the fixed dimension
of $S_i$ gives (SS.8).  Equivalently write each fixed dimension
as $b$ minus its moving rank; the $b\,\operatorname{ind}(x_i)$
term then follows directly.

Above the $i$th branch value of $C\to\mathbf P^1$, the points
of $X$ correspond to the cycles of $x_i$, and their inertia
operators in $\tau$ are conjugate to $S_i^{|c|}$.
This proves the formula for $D'$ even when some of these
operators are the identity and are omitted as genuine marks.
Riemann--Hurwitz for $X\to\mathbf P^1$ gives
$\sum_i\operatorname{ind}(x_i)=2n-2+2h$.
For a nontrivial irreducible representation $T$ of a finite
sphere cover, the usual Euler formula is
$\mu_T=-2\dim T+\sum_i\operatorname{rank}(1-T(g_i))$.
Sum (SS.8), use $\mu_V=6$, and use this formula for $U$ to obtain
(SS.9).  The version for the trivial representation has its
$2\dim T^J=2$ invariant correction.  Equivalently the general
Euler formula over $X$ contains $2\dim\tau^H$; (SS.2) and
Frobenius show $\dim\tau^H=\dim U^J$, since $V$ is nontrivial.
\end{proof}

\paragraph{A genus bound which retains the extra summand.}
The accepted weak all-target standard estimate
$\nu(W(x))\ge (n-2)\operatorname{ind}(x)/n$, together with
spectral convolution with $U$, gives
\[
 2(n-1)b+6
  \ge\frac{b(n-2)}n\bigl(2n-2+2h\bigr),\qquad
 6\ge\frac{2b}{n}\bigl((n-2)h-2(n-1)\bigr).
 \tag{SS.15}
\]
Both inequalities concern the actual natural quotient and require
no lower support bound for the partner.  If $h\ge3$, they imply
$b\le3n/(n-4)$.  Consequently $b\ge4,n\ge17$ force $h\le2$.
They do not force $h\le1$.

If the $V$-period rank is nine and $h\ge1$, the marked-base
dimension for the connected $H/N$-cover $Y\to X$ is at most
$3h-3+D'$ (with the usual $h=1$ convention, equal to the
number of genuine marks when that number is positive).
Together with (SS.9) this yields the necessary condition
\[
               \mu_U\ge (2b-3)(h-1)+3.
 \tag{SS.16}
\]
For an unramified elliptic base the moduli dimension is only one,
so rank nine is already impossible; no exceptional convention
weakens this necessary condition.  In particular $h=1$ requires
$\mu_U\ge3$, and $h=2$ requires $\mu_U\ge2b$.
These are bounds on a surviving summand, not a proof of its
vanishing.  For $h=0$ the analogous necessary bound is
$\mu_U\ge6-2b$.

There is also a useful upper bound.  Every nonidentity target
eigenvalue on $W(x)$ has multiplicity at most $n/2$: each
supporting permutation cycle has length at least two.
The non-one eigenspaces of $U(g_i)$ therefore contribute at least
$(n-2)/2$ each to the moving rank of $V(g_i)$, counted with their
multiplicities.  Summing and using the nontrivial-$U$ Euler formula
gives
\[
  \frac{n-2}{2}(\mu_U+2b)\le2(n-1)b+6,
  \qquad
  \mu_U\le 2b+\frac{4b+12}{n-2}.
 \tag{SS.17}
\]
For example, when $h=2$ and $n>4b+14$, the two integer bounds
force $\mu_U=2b$.  No generic-endomorphism conclusion for the
resulting genus-two cover is assumed.

\paragraph{What follows if the extra summand vanishes.}
Suppose $\mu_U=0$ and the original exact-$K$ sixfold period map
has rank nine.  All Galois conjugate $U$-summands also have zero
cohomology, since $H^1(C,\mathbf Q)$ is rational.
Equation (SS.7) identifies the complete complex $\tau$-period
variation with the original one, so its rank is nine.
For $b\ge2$, (SS.9) forces $h=0$: if $h=1$ then $D'=6$ and
the marked elliptic-base moduli has dimension at most six;
if $h=2$, only $b=2,3$ are possible, and the moduli dimensions
are at most five and three; $h\ge3$ is impossible.
Finite monodromy choices over a fixed marked curve are discrete,
so these marked-base dimensions bound the period rank.
Thus $Y\to X$ supplies a lower-degree, degree-$b$ complex
multiplicity variation over a rational base with at least twelve
genuine marks.  Its moving-rank total is $2b+6$, so necessarily
$b\ge3$.

There is also a precise rational statement.  Let $B$ be the
original primitive $K$-factor and let $A$ be a primitive $L$-factor
for $\tau$ in $J(Y)$, using the models constructed above.  Put
$\ell=[L:K]$.  Upon inducing the rational representation
underlying the $L$-model of $\tau$, (SS.2) and all its conjugates
give the rational $U$-orbit once and the rational $V$-orbit
$\ell$ times.  The first has zero Hom into $H^1(C)$.
Rational Frobenius therefore gives
\[
                              A\sim B^\ell
 \tag{SS.10}
\]
in families.  These are rational deck correspondences and the
actual pullback along $C\to Y$; the actual Jacobian polarization
on $A$ is retained.  No equality of integral polarization types
or arbitrary choice of a polarization is inferred.

This is a descent with an inactive extra constituent.  Even when
$\mu_U=0$, the representation $\operatorname{Ind}_H^J\tau$
is still $U\oplus V$ and is not irreducible.  Thus this observation
alone does not supply an irreducible-induction presentation of
$V$ of smaller inducing degree.  It also does not prove
$\mu_U=0$.  For $b=1$, the exact scalar-standard situation is
already treated by D-BI-08 and is not reclassified here.

\paragraph{The Hodge dimensions when the summand survives.}
Write $(p,q)$ for the Hodge multiplicities of $U$ at the chosen
embedding.  Then (SS.7) gives the exact multiplicities
$(p+3,q+3)$ for $\tau$, with $p+q=\mu_U$.
The two summands are fixed subvariations.  For the actual invariant
pairing their cross terms vanish by the different $J$-isotypic
types; on the multiplicity forms the two restrictions may carry
different fixed positive scalar factors.
Consequently the period differential of (SS.7) is the direct sum
of the two period differentials, and its kernel is their
intersection.  Its rank is at least the rank of the $V$-period
map, not at most that rank.  The block-diagonal period domain
has dimension $pq+9$, which supplies no bound below nine.
Even when $p=0$ or $q=0$, the $U$-summand is constant at that
embedding and the $V$-summand can still have rank nine at the
level of polarized Hodge variations.

For clarity this last logical limitation has an abstract exact-$K$
example: take the rank-nine tautological variation on a small open
set in the $U(3,3)$ period domain and form its direct sum with a
fixed polarized $K$-Hodge structure of signature $(1,0)$.
The resulting signatures are $(4,3)$ and the period rank is
still nine, with nonzero extra summand.  This is only an example
of polarized Hodge variations satisfying the displayed direct-sum
identity.  It is not a finite-group curve construction or a
realization of the tensor representation (SS.1).  Any further
vanishing or rank conclusion must use those additional geometric
or representation-theoretic conditions.

\subsection*{HD-SS03. An actual Euler-only counterexample}

\begin{proposition}[HD-SS03]
For every even $n=2m\ge6$, there is a connected genus-zero-base
finite-group Hurwitz family with twelve genuine inertia entries,
faithful irreducible $V=\operatorname{Std}_n\otimes U$, and
$\mu_V=6$, but $\mu_U=2$.  Its exact coefficient field is
$\mathbf Q$, and its $V$-period rank is at most six.
\end{proposition}

\begin{proof}
Take $J=C_2\times S_n$, let $U=\chi$ be the nontrivial character
of the first factor, and set
\[
 a_0=(1\ 2)(3\ 4)\cdots(n-1\ n),\qquad
 b_0=(2\ 3)(4\ 5)\cdots(n\ 1),
 \qquad t_1=(1\ 2),\quad t_2=(1\ 3).
 \tag{SS.11}
\]
Use the tuple
\[
 (-1,a_0),(-1,b_0),(-1,b_0),(-1,a_0),
 (1,t_1),(1,t_1),(1,t_2),(1,t_2),
 (1,t_1),(1,t_1),(1,t_2),(1,t_2).
 \tag{SS.12}
\]
It is product one and all entries are genuine involutions.
The powers of $a_0b_0$ conjugate $t_2$ to the edges of a connected
cycle on the odd letters; conjugation by $a_0$ supplies a connected
cycle on the even letters, and $t_1$ joins the two.
Transpositions on a connected graph generate $S_n$.
The full generated subgroup also projects onto $C_2$.
If its kernel over $S_n$ were trivial, it would be the graph of
a homomorphism $S_n\to C_2$.  The entry $(1,t_1)$ would force
that homomorphism to be trivial, since transpositions are
conjugate and generate $S_n$, contrary to the first entry.
Hence the tuple generates $J$.

The deleted permutation representation is irreducible and has
no nonidentity scalar permutation for $n\ge3$.  Therefore
$\chi\otimes\operatorname{Std}_n$ is faithful and irreducible,
and has exact character field $\mathbf Q$.
At each of the first four entries its moving rank is $m-1$;
at each of the last eight it is one.  Thus
\[
 \sum_i\operatorname{rank}(1-V(g_i))
   =4(m-1)+8=2(n-1)+6.
 \tag{SS.13}
\]
The product-one generation criterion, in the accepted D-NR
Hurwitz framework, realizes an actual connected cover and its
Hurwitz component.  Euler gives $\mu_V=6$.  The representation
has a real rational orthogonal model, so complex conjugation
interchanges its two Hodge multiplicities; these are $(3,3)$.

The quotient for $\chi$ is a double cover of $\mathbf P^1$
branched at exactly four points, hence is elliptic and
$\mu_\chi=2$.  The natural $n$-sheeted quotient $X=C/H$ has
total index $4m+8=2n+8$, so $g(X)=5$.
The induced double cover $Y\to X$ is connected and unramified:
at a negative-character branch all permutation cycles have
length two, and at every other branch the character is trivial.
Consequently
\[
 g(Y)=9,\qquad
 \dim H^1(Y,\mathbf C)^{-}=2(9-5)=8=2+6.
 \tag{SS.14}
\]
This is exactly the extra summand in (SS.7).
The primitive multiplicity factor for $V$ is an abelian
threefold, whose Siegel period domain has dimension
$3(3+1)/2=6$.  Thus this example cannot have period rank nine
and is outside the exact imaginary-quadratic target.
\end{proof}

\paragraph{The surviving question.}
Neither the Euler budget nor twelve genuine marks alone proves
$\mu_U=0$.  HD-SS03 refutes only that shortcut.
In the original exact-$K$ setting, HD-SS01 forces $U$ to be
nonselfdual, which the counterexample does not satisfy.
This section proves neither that $\mu_U>0$ is impossible in
that setting nor that its $V$-period rank is below nine.
It retains the exact extra summand, the proper-field possibility
$L\supset K$, and the difference between an inactive-summand
descent and irreducible induction.
